\documentclass[11pt]{article}
\usepackage[margin=0.75in]{geometry}
\usepackage{fontspec}
\usepackage{setspace}
\usepackage{hyperref}
\hypersetup{colorlinks=true, urlcolor=blue, linkcolor=black}
\pagestyle{empty}
\setlength{\parskip}{0.75em}
\setlength{\parindent}{0pt}
\setstretch{0.94}
\setmainfont{Times New Roman}

\begin{document}

\begin{flushleft}
Ke Cheng HSU (Clio) \\
Department of Economics, The Hong Kong University of Science and Technology \\
Clear Water Bay, Kowloon, Hong Kong $\cdot$ +852 6840 9753 \\
\href{mailto:khsu@connect.ust.hk}{khsu@connect.ust.hk} $\cdot$ \href{https://cliohsu1997.github.io/}{cliohsu1997.github.io}
\end{flushleft}

7 November 2026

\vspace{0.15em}
Faculty Search Committee \\
Zhejiang University International Business School \\
Haining, China

\vspace{0.15em}
Dear Members of the Search Committee,

I am writing to apply for the Postdoctoral Position in Marketing at the Zhejiang University International Business School. I am a Ph.D.\ candidate in Economics at the Hong Kong University of Science and Technology. I expect to graduate in June 2027 and can start afterward. My fields are Empirical Industrial Organization and Digital Economics. Within the areas listed in the posting, my work sits in digital marketing, organizations and strategy in the digital age, and information systems and business analytics.

My job market paper, ``Designing Return Policy in E-Commerce: A Trade-off between Purchase and Return Efficiency,'' studies a core digital-marketing decision: how an online seller should design the return offer. A stricter policy curbs inefficient returns after delivery, but the lower return rate makes a higher product price more attractive for the seller and can exclude efficient purchases. Using transaction data from a Taobao apparel retailer, I estimate a structural model of purchase, return, and whether to buy return-shipping insurance under monopoly pricing across six categories. Automatic coverage of return shipping performs better when the margin, net scrap value, and uncertainty are high, because it expands efficient purchases. Optional coverage performs better when those attributes are low, because it curbs excessive returns. A menu that pairs a discounted price with a medium restocking fee balances the two margins and performs better than both insurance policies for most categories.

A second working paper, ``Product Design with Heterogeneous Search Costs'' (with Pak Hung Au), studies how firms choose product design when consumers differ in search costs. Search costs shape which products consumers search. Consumers with low search costs seek niche products, and consumers with high search costs seek standardized products. A fall in search costs can raise prices when this shift in which product consumers search outweighs the effect that encourages more search. In two ongoing projects, with Kohei Kawaguchi and Kei Kawai, I use Japanese bookstore data to study how resale price maintenance encourages retailers to hold inventory in low-demand states, yet leaves unsold copies that cannot be cleared by cutting the price and are returned. With Ce Bian and Liang Zhao, I build a theoretical model of pass-through when products can be returned, and show that a tax can lower the equilibrium price and raise welfare by reducing the return rate.

At HKUST I was a teaching assistant from 2020 to 2025 for undergraduate and Ph.D.\ courses, with an average evaluation of 4.72/5, and in 2026 I designed and lectured the mathematics camp for the MSc in Economics, receiving 4.53/5. I attach my CV and the two papers named above. Pak Hung Au, Kohei Kawaguchi, and Kei Kawai will provide recommendation letters. You can reach me at \href{mailto:khsu@connect.ust.hk}{khsu@connect.ust.hk} or +852 6840 9753.

\vspace{0.35em}
Sincerely,

\vspace{0.35em}
Ke Cheng HSU \\
Ph.D.\ Candidate, Department of Economics, HKUST

\end{document}
